\label{sec:terms}

В этом разделе зафиксированы термины и обозначения, используемые далее в
статье. Приведённые пояснения служат кратким справочником; формальные
определения и допущения вводятся в следующем разделе.

\subsection{Core Terms}

\begin{description}[style=nextline]
  \item[Единица работы] Условный объём работы, трудоёмкость которого без ИИ
  равна $X$. В иллюстративном сценарии единицы работы связываются с оценкой в
  story points, однако модель не требует конкретной шкалы сложности.

  \item[Задача] Неделимая на выбранном уровне декомпозиции часть проекта.
  Задача имеет сложность $Z_i$, режим выполнения $r$ и нормированную длительность
  $k_i^{(r)}$. При изменении уровня декомпозиции одна задача может быть
  заменена несколькими подзадачами.

  \item[Трудоёмкость] Объём работы, выраженный во времени выполнения без ИИ.
  В исходном сценарии с одним разработчиком трудоёмкость совпадает с
  длительностью; при параллельном выполнении суммарная трудоёмкость и
  календарный срок проекта различаются.

  \item[Календарное и активное время] Календарное время измеряется от начала
  до завершения задачи и включает ожидание. Активное время учитывает только
  интервалы непосредственной работы исполнителя. Эти показатели, как и
  пропускная способность потока задач, не являются взаимозаменяемыми.

  \item[Принятый результат] Результат, удовлетворяющий заданному критерию
  приёмки. Время проверки, исправления и повторной проверки включается в
  длительность задачи; простое завершение агентной генерации не считается
  завершением работы.

  \item[Исходный сценарий] Последовательное выполнение заданного объёма работ
  одним разработчиком без ИИ. Его продолжительность обозначается $T_h$ и
  используется как база для расчёта ускорения. В англоязычной литературе
  употребляется термин \emph{baseline}.

  \item[Нормированная длительность с ИИ] Отношение длительности задачи с ИИ к
  базовой оценке $Z_iX$. Обозначается $k_i^{(r)}$ и зависит от задачи,
  инструмента и режима. Это описательная величина, а не универсальная
  характеристика модели ИИ или причинная оценка без отдельного дизайна.

  \item[ИИ-агент для программирования] Программная система на основе модели
  ИИ, способная получать задачу, работать с кодом и инструментами и возвращать
  результат для автоматической либо человеческой приёмки. Краткая форма в
  статье --- \emph{агент}.

  \item[Человеко-агентная ячейка] Рассматриваемая в статье организационная
  единица: один разработчик ставит задачи одному или нескольким агентам и
  принимает их результаты. Модель нескольких взаимодействующих ячеек остаётся
  направлением дальнейшего развития.

  \item[Режим работы] Способ взаимодействия человека и агента при выполнении
  задачи. Различаются интерактивный, делегированный и полностью автономный
  режимы. Режим определяет не продукт сам по себе, а порядок постановки,
  выполнения, проверки и исправления задачи.

  \item[Интерактивный режим] Разработчик и агент совместно выполняют одну
  задачу, последовательно обмениваясь запросами и результатами. В предельном
  случае человеческая составляющая занимает всё время задачи: $h_i=k_i$.

  \item[Делегированный режим] Разработчик ставит задачу, агент выполняет её
  автономно, после чего человек проверяет результат, принимает его или
  возвращает на доработку.

  \item[Полностью автономный режим] Агент самостоятельно обнаруживает задачу,
  выполняет её и подтверждает результат без обязательного участия человека.
  Этот режим описывает границу применимости модели и подробно не исследуется.
  Его не следует смешивать с автономной agent-фазой делегированной задачи:
  такая фаза освобождает разработчика на время исполнения, но постановка и
  приёмка результата по-прежнему требуют human-фаз.

  \item[Конфигурация процесса] Полное описание организации работы: настроенный
  параллелизм $P$, назначение режимов задачам, используемый инструмент и
  протокол постановки и приёмки. При фиксированных инструменте и протоколе
  применяется сокращённая запись $\sigma=(P,r)$ для единого режима либо
  $\sigma=(P,\rho)$ для разных режимов. Сравнивать следует конфигурации целиком,
  поскольку изменение любого из этих элементов может изменить коэффициенты
  $k_i^{(r)}$, фазовые профили и издержки.

  \item[Декомпозиция] Разбиение проекта на задачи и подзадачи с явными
  границами и зависимостями. Декомпозиция определяет доступный параллелизм и
  состав критического пути; чрезмерное дробление может увеличить издержки
  постановки и приёмки.

  \item[Параллелизм] Число одновременно доступных потоков выполнения с ИИ.
  Обозначается $P$. В задачах составления расписания $P$ равно целому числу
  исполнителей, а в аналитических оценках может трактоваться как эффективное,
  усреднённое за период значение. Число запущенных агентов не обязательно
  совпадает с фактически достижимым параллелизмом.

  \item[Граф зависимостей] Ориентированный ациклический граф $G=(V,E)$,
  вершины которого соответствуют задачам, а дуги задают отношения
  предшествования. Англоязычное сокращение --- DAG
  (\emph{directed acyclic graph}).

  \item[Расписание] Назначение задач исполнителям и моментов их начала и
  завершения. Расписание называется допустимым, если соблюдает отношения
  предшествования, порядок фаз, число агентных потоков, локальное закрепление
  агента за начатой задачей и ограничение человеческого ресурса.

  \item[Фаза задачи] Непрерывный этап выполнения задачи с однозначно заданным
  типом ресурса. В модели различаются автономные агентные и человеческие фазы;
  последние требуют внимания разработчика и одновременно сохраняют
  закрепление соответствующего агентного потока.

  \item[Критический путь] Самый длинный по времени путь в графе зависимостей.
  Его длина $L$ задаёт нижнюю границу срока проекта: увеличение числа агентов
  не сокращает задачи, последовательно связанные на этом пути.

  \item[Время завершения всех работ] Время от начала проекта до завершения
  последней задачи с принятым результатом. Это основной временной показатель
  статьи. В теории расписаний ему соответствует термин \emph{makespan}.

  \item[Суммарная работа] Сумма длительностей всех задач с ИИ при выполнении в
  отдельных потоках. Обозначается $W$. Величина $W/P$ задаёт идеализированную
  оценку при равномерном распределении нагрузки.

  \item[Агентная составляющая] Интервалы выполнения задачи, в течение которых
  агент работает автономно, а разработчик свободен. Нормированная
  продолжительность обозначается $a_i^{(r)}$.

  \item[Человеческая составляющая] Интервалы, в течение
  которых разработчик занят постановкой, чтением, проверкой или направлением
  исправлений по задаче. Нормированная продолжительность обозначается
  $h_i^{(r)}$.

  \item[Человеческий ресурс] Время и внимание одного разработчика, общие для
  всех потоков. Ресурс имеет мощность 1: человеческие интервалы разных задач
  не могут пересекаться.

  \item[Координационные издержки] Дополнительные затраты, возникающие при
  параллельной работе нескольких потоков. Межагентные интеграционные издержки
  учитываются функцией $C(P)$, а затраты разработчика на переключение между
  потоками --- множителем $\gamma(P)$.

  \item[Локальная блокировка агентного потока] Правило уровня M4, согласно
  которому начатая задача удерживает назначенного агента до принятия результата.
  Если задача запросила внимание занятого разработчика, агент прекращает работу
  и остаётся занятым ожиданием; остальные агенты продолжают свои задачи.
  Блокировка одного потока не означает остановки всей системы.

  \item[Очередь к человеческому ресурсу] Множество запросов агентов, ожидающих
  обслуживания одним разработчиком. Ожидание не входит в полезное человеческое
  время $H$, но увеличивает занятость закреплённых агентов и может привести к
  превышению фактического времени над нижней оценкой $B_4$.

  \item[Переиспользование заблокированного потока] Запуск другой задачи на месте
  агента, ожидающего человека. Такая организация, как и запуск приоритетного
  дополнительного потока сверх $P$, в статье не рассматривается.

  \item[Ускорение] Отношение времени исходного сценария к времени выбранной
  конфигурации с ИИ. Значение больше единицы означает сокращение срока, единица
  --- отсутствие временного эффекта, значение меньше единицы --- замедление.

  \item[Нижняя оценка времени] Значение, меньше которого время завершения быть
  не может. На уровнях M1--M4 нижняя оценка не обязательно достижима даже при
  оптимальном расписании.

  \item[Узкое место] Задача, путь или ресурс, который в рассматриваемой
  конфигурации определяет нижнюю оценку времени. В модели узким местом может
  быть суммарная нагрузка $W/P$, самая длинная задача $M_N$, критический путь
  $L$ или человеческий ресурс $\gamma(P)H$.

  \item[Уровни M0--M4] Последовательные уточнения модели. Уровень M0
  предполагает идеальную делимость работы; M1 учитывает неделимость задач, M2
  --- граф зависимостей, M3 переопределяет длительности с учётом
  координационных издержек, а M4 вводит фазовый порядок, ограниченный
  человеческий ресурс и локальную блокировку.
\end{description}

\subsection{Notation}

\begin{description}[style=nextline]
  \item[$X$] Базовая трудоёмкость единицы работы без ИИ.

  \item[$N$] Число задач в проекте; $i\in\{1,\dots,N\}$ --- индекс задачи.

  \item[$Z_i$] Сложность $i$-й задачи в единицах $X$.

  \item[$R$, $r$] Множество допустимых режимов работы и выбранный режим.

  \item[$\rho$] Отображение, назначающее каждой задаче собственный режим:
  $\rho\colon\{1,\dots,N\}\to R$.

  \item[$T_i^{ai}(r)$] Наблюдаемое время выполнения $i$-й задачи с ИИ в одном
  потоке режима $r$ при полной доступности разработчика.

  \item[$k_i^{(r)}$] Нормированная длительность $i$-й задачи с ИИ в режиме $r$:
  $k_i^{(r)}=T_i^{ai}(r)/(Z_iX)$. Значения меньше, равные и больше единицы
  соответствуют меньшей, равной и большей длительности относительно базы.

  \item[$\hat{k}_i^{(r)}$] Прогнозная оценка $k_i^{(r)}$, полученная до начала
  проекта по сопоставимым задачам и режимам.

  \item[$a_i^{(r)}$, $h_i^{(r)}$] Агентная и человеческая составляющие
  коэффициента $k_i^{(r)}$, причём $k_i^{(r)}=a_i^{(r)}+h_i^{(r)}$.

  \item[$P$, $P_h$] Параллелизм процесса с ИИ и параллелизм исходного
  сценария. В статье принимается $P_h=1$.

  \item[$\sigma$] Сокращённая запись конфигурации процесса при фиксированных
  инструменте и протоколе приёмки: $\sigma=(P,r)$ или $\sigma=(P,\rho)$.

  \item[$T_h$] Время последовательного выполнения всего проекта одним
  разработчиком без ИИ.

  \item[$T_{ai}$, $T_{ai}^{(0)}$] Время выполнения проекта с ИИ и его
  идеализированное значение на уровне M0.

  \item[$W$] Суммарная работа с ИИ:
  $W=X\sum_{i=1}^{N}Z_i k_i$.

  \item[$w_i$] Вес, или длительность, отдельной задачи с ИИ:
  $w_i=Z_i k_i X$.

  \item[$A$] Суммарная автономная агентная работа:
  $A=X\sum_i Z_i a_i$.

  \item[$M_N$] Длительность самой длинной задачи:
  $M_N=X\max_i(Z_i k_i)$.

  \item[$G=(V,E)$] Ориентированный ациклический граф зависимостей между
  задачами.

  \item[$L$] Длина критического пути в $G$ с весами $w_i$.

  \item[$W_3(P)$, $L_3(P)$] Суммарная занятость агентных потоков и длина
  критического пути на уровне M3 после применения $C(P)$ только к автономным
  агентным фазам.

  \item[$W_4(P)$, $L_4(P)$] Аналогичные величины уровня M4 после дополнительного
  учёта множителя $\gamma(P)$ в длительности человеческих фаз.

  \item[$C(P)$] Множитель межагентных координационных и интеграционных
  издержек; $C(1)=1$.

  \item[$H$] Суммарное человеческое время:
  $H=X\sum_{i=1}^{N}Z_i h_i^{(r)}$.

  \item[$\gamma(P)$] Множитель затрат на переключение внимания разработчика
  между параллельными потоками; $\gamma(1)=1$.

  \item[$Q(\pi)$] Суммарное время блокировки назначенных агентов в очереди к
  разработчику в конкретном расписании $\pi$. Это эндогенная характеристика
  расписания, не входящая в $k_i$ и $H$.

  \item[$T(\pi)$, $T^{*}$] Время завершения конкретного допустимого расписания
  $\pi$ и минимальное время среди всех допустимых расписаний соответственно.

  \item[$\bar{k}_w$] Средневзвешенная нормированная длительность:
  $\bar{k}_w=\sum_i Z_i k_i/\sum_i Z_i$.

  \item[$\bar{a}_w$] Средневзвешенная нормированная агентная работа:
  $\bar{a}_w=\sum_i Z_i a_i^{(r)}/\sum_i Z_i$.

  \item[$\bar{h}_w$] Средневзвешенная нормированная человеческая работа:
  $\bar{h}_w=\sum_i Z_i h_i^{(r)}/\sum_i Z_i$.

  \item[$B_0,\dots,B_4$] Нижние оценки времени завершения в последовательных
  уточнениях M0--M4. Переход между уровнями может добавлять ресурсное
  ограничение и переопределять эффективные длительности.

  \item[$S_E$, $S_E^{\mathrm{ach}}$] Идеализированный и достижимый
  коэффициенты ускорения.

  \item[$S_{\mathrm{target}}$] Целевой коэффициент ускорения.

  \item[$\mu$] Отношение длительности узкого места к исходному времени $T_h$:
  $\mu=M_N/T_h$ на уровне M1 и $\mu=L/T_h$ на уровне M2. Символ $m$ в
  вычислительных экспериментах зарезервирован для числа частей декомпозиции.

  \item[$\nu_i$] Математическое ожидание случайной нормированной длительности:
  $\nu_i=\mathbb{E}[k_i]$.

  \item[$b$, $b_0$] Отношение длительности узкого места к средней нагрузке
  $W/P$ и заданная верхняя граница этого отношения, используемая как правило
  декомпозиции.

  \item[$\kappa$] Общий положительный множитель, применяемый к длительностям
  всех agent- и human-фаз при анализе масштабной инвариантности. Отдельного
  масштабирования только суммарных $k_i$ для вывода уровня M4 недостаточно.
\end{description}
